% ===========================================================================
%  sviluppatore/05-estendere/05-regime.tex
% ===========================================================================
\section{Un regime: due file, in quest'ordine}\label{sez:nuovo-regime}

\pkgRegime{} è l'unica famiglia divisa in due artefatti distinti, ed è la
divisione giusta: \emph{una regola serve più mercati}. Azioni e cripto usano
entrambe \cod{trend\_vol}, con parametri diversi.

\figuraTikz{catena-regime}
  {Dal simbolo allo stato. I due riquadri con il file sotto sono le due cose
   che si scrivono.}
  {catena-regime}

\subsection{La regola}

\begin{codiceInterno}{python}{Il contratto della regola}{regime/logics/\_\_init\_\_.py:7}
    META = {
        "key": "mia_regola",
        "label": "La mia regola",
        "description": "a cosa serve",
        "defaults": {"finestra": 60},
    }

    def compute(close: pd.Series, p: dict) -> pd.Series:
        ...   # Series di 'bull' / 'neutral' / 'bear'

Due vincoli non negoziabili:

1. **Causalita'**: il valore al tempo t usa solo dati fino a t. Niente
   statistiche sull'intero campione (quantili full-sample = look-ahead).
2. **Etichette**: sempre ``bull``/``neutral``/``bear``, cosi' i
   ``regime_weights`` gia' dichiarati dalle strategie continuano a valere.
\end{codiceInterno}

\noindent
Gli helper causali condivisi stanno in \file{regime/logics/\_base.py}:
\simbolo{\_realized\_vol}, \simbolo{\_expanding\_median},
\simbolo{\_hysteresis}. Vanno usati: riscrivere un quantile a mano è il modo
più comune di introdurre una statistica full-sample senza volerlo.

Le tre regole installate:

\begin{center}
\begin{tabular}{@{}L{26mm}L{54mm}L{58mm}@{}}
\toprule
\textbf{Chiave} & \textbf{Che cosa misura} & \textbf{Che cosa significano le
  etichette} \\
\midrule
\rigaTenue
\cod{trend\_vol} & Posizione rispetto alla media mobile lunga, corretta per
  volatilità realizzata & bull = sopra il trend e volatilità accettabile;
  bear = sotto il trend, o volatilità alta \\
\cod{vol\_state} & Volatilità realizzata contro i quantili espansivi della
  propria storia & bull = calmo; bear = sotto stress. Per attività senza
  deriva \\
\rigaTenue
\cod{trend\_range} & Rapporto di varianza sui rendimenti: la persistenza, non
  la direzione & bull = in tendenza; bear = laterale. Indipendente dalla
  deriva \\
\bottomrule
\end{tabular}
\end{center}

\subsection{Il mercato}

\begin{codiceInterno}{python}{Il contratto del mercato}{regime/markets/\_\_init\_\_.py:12}
    META = {
        "key": "equity",
        "label": "Azioni",
        "logic": "trend_vol",          # regola ereditata da regime/logics.py
        "driver": ["SPY"],             # serie che GUIDA il regime (in ordine)
        "match": [r"^[A-Z]{1,5}$"],    # simboli assegnati a questo mercato
        "priority": 50,                # match piu' bassa = valutata prima
        "enabled": True,
        "params": {"trend_window": 200},   # override sui default della logica
    }
\end{codiceInterno}

\noindent
Un esempio reale, che mostra come si usa \cod{priority} e che cosa
significa \cod{"self"} fra i driver:

\begin{codiceInterno}{python}{Il mercato valutato per primo}{regime/markets/forex.py:17}
META = {
    "key": "forex",
    "label": "Forex",
    "logic": "vol_state",
    "driver": ["DXY", "self"],
    "match": [r"^(EUR|USD|GBP|JPY|CHF|AUD|NZD|CAD|SEK|NOK)"
              r"(EUR|USD|GBP|JPY|CHF|AUD|NZD|CAD|SEK|NOK)$",
              r"^[A-Z]{3}/[A-Z]{3}$",
              r"^DXY$"],
    "priority": 10,          # valutato per primo: pattern molto specifico
    "enabled": True,
    "params": {
        "vol_window": 21,
        "min_obs": 252,
        # banda ampia + isteresi lunga: un regime deve essere uno STATO, non
        # una reazione. Valori strutturali, non tarati sui dati.
        "calm_quantile": 0.30,      # sotto questo quantile: risk-on
        "stress_quantile": 0.80,    # sopra questo quantile: risk-off
        "hysteresis_days": 10,
    },
}
\end{codiceInterno}

\begin{nota}{\cod{driver} è una lista di ripieghi}
Si valutano in ordine: il primo disponibile vince. \cod{"self"} significa «il
simbolo stesso». Se nessuno dei driver è disponibile il mercato non modula
--- cioè si comporta come neutro --- e lo dichiara in un avviso nel registro,
non in un errore. Un mercato che non modula non si distingue da un mercato
che modula con pesi \num{1}, se non leggendo il registro.
\end{nota}

\subsection{\cod{META} di soli letterali: qui è un vincolo vero}

Per le regole di regime il sistema legge il \cod{META} con
\simbolo{ast.literal\_eval}, cioè \textbf{senza eseguire il file}. La
conseguenza è netta: un valore come \cod{"x": np.pi} non è un letterale, la
lettura fallisce, e la regola diventa invisibile.

\begin{attenzione}{La \cod{key} deve coincidere con il nome del file}
Per le regole di regime il controllo è esplicito ed è l'unico caso in cui il
disallineamento viene rifiutato. Negli altri casi non lo è, e il plugin
compare con un nome e si cerca con un altro.
\end{attenzione}

\subsection{Il merge dei parametri}

I \cod{params} dichiarati dal mercato scavalcano i \cod{defaults} dichiarati
dalla regola, e il punto in cui accade è uno solo:

\begin{codiceInterno}{python}{Dove i due artefatti si incontrano}{regime/logics/\_\_init\_\_.py:89}
def apply(logic_key: str, close, params: dict | None = None):
    if logic_key not in REGISTRY:
        raise KeyError(f"unknown regime rule: '{logic_key}' "
                       f"(available: {sorted(REGISTRY)})")
    entry = REGISTRY[logic_key]
    p = dict(entry["defaults"])
    p.update(params or {})
    return entry["fn"](close, p)
\end{codiceInterno}

\begin{pratica}{Aggiungere un regime}
\begin{enumerate}
  \item \file{regime/logics/mia\_regola.py}: \cod{META} di soli letterali con
        \cod{key} uguale al nome del file, e \cod{compute(close, p)} che
        restituisce \cod{'bull'}/\cod{'neutral'}/\cod{'bear'} allineati a
        \cod{close}, causali, con isteresi.
  \item \file{regime/markets/mio\_mercato.py}: \cod{META} con almeno
        \cod{key} e \cod{logic}; la regola nominata deve \emph{già} esistere.
  \item Scaricare le serie dei driver:
        \file{python scripts/download\_regime\_drivers.py}.
  \item Scorciatoia: \cod{regime.logics.scaffold(key)} genera lo scheletro
        corretto, ed è quello che usa l'editor dell'interfaccia.
\end{enumerate}
Le guide sono \file{regime/GUIDE\_REGIMI.md} («come costruirne uno») e
\file{strategies/GUIDE\_REGIME.md} («come usarlo dentro una strategia»): sono
due documenti distinti e si rimandano a vicenda.
\end{pratica}
